\section{第\ref{sec:glutcomputer}章　\nt{GLUT}ニューロンは何を計算するか}

この章の論理的な主張は、非負の重みからなる回路についての定理であり、章の中で導いたものである。実験が裏づけるのは、ニューロンのモデルと、これらの定理に従って組んだ回路（同時検出器、遅延線、自己維持する群、連鎖）が脳に実際に存在することである。

{\footnotesize
\setlength{\tabcolsep}{3pt}\renewcommand{\arraystretch}{1.15}
\begin{longtable}{>{\raggedright}p{0.5cm}>{\raggedright}p{4.0cm}>{\raggedright}p{5.6cm}>{\raggedright}p{2.3cm}>{\raggedright\arraybackslash}p{1.7cm}}
\toprule
No. & 主張 & 実験と測定されたもの & 出典 & 状態 \\
\midrule\endfirsthead
\toprule
No. & 主張 & 実験と測定されたもの & 出典 & 状態 \\
\midrule\endhead
1 & ニューロンは閾値とリセットをもつRC回路である~\eqref{eq:glutneuron} & ラット皮質の錐体ニューロン（スライス）の細胞体に、生きた脳のシナプス入力の流れに似た雑音電流を注入した。平均発火率の電流の平均とばらつきへの依存性は、遅い発火率適応を加えれば漏れ積分発火ニューロンで記述できる。$\tau$と遅延の範囲は実験を引用せずに章で示した & \cite{Rauch2003} & 測定済み \\
2 & モデル~\eqref{eq:glutneuron}は単純化であり、入力の枝そのものが局所スパイクを出す & マウス視覚野の錐体ニューロンの樹状突起からin vivoで直接記録：視覚刺激がNMDA受容体に依存する局所的な樹状突起スパイクを起こし、それを抑えるとニューロンの方位選択性が広がる & \cite{Smith2013} & 測定済み \\
3 & 同時性の窓$\Delta_{\max}=\tau\ln\frac{w}{\theta-w}$~\eqref{eq:window}。$\theta=1.5w$では$0.7\tau$ & モデル~\eqref{eq:glutneuron}からの帰結。狭い加算窓の直接測定は速い抑制をもつニューロンについてある（第\ref{sec:gaba}章、その表の3行目） & 章内の導出 & 理論 \\
4 & \nt{GLUT}ニューロンだけのネットワークは入力の単調関数しか計算しない（命題~\ref{prop:monotone}） & 章内の証明：非減少の右辺による静寂からの反復。これはモデルについての主張であり、抑制のない生きたネットワークで検証した実験はない & 章内の導出 & 理論 \\
5 & 感覚器は配線ペアを与える。刺激の増加に応答する細胞と減少に応答する細胞がある & 網膜：双極細胞の段階ですでに錐体の信号はオンチャネルとオフチャネルに分かれる。サルでオン双極細胞を薬理学的に遮断：両チャネルは皮質まで別々に進み、遮断後は光の増加の検出が悪化するが、減少の検出は悪化しない & \cite{Schiller1992} & 測定済み \\
6 & 2線式符号があれば、\nt{GLUT}ニューロンのネットワークは任意の論理関数を共通のクロックに頼らず非同期に計算でき、代償はニューロン数が2倍になること（命題~\ref{prop:dualrail}、~\eqref{eq:dualrail}、図~\ref{fig:dualxor}） & 2線式符号と信号自体による完了検出は、遅延に依存しない非同期回路の標準的な技法である。ニューロン6個の排他的OR回路は章の中で組んだ。脳がまさに2線式符号で論理関数を計算していることを示す実験はない & \cite{Sparso2001} & 理論 \\
7 & ヒトで両耳への音の到着時刻の差は最大$\approx0.6\ms$であり、脳は約10マイクロ秒を区別する & 二肢選択実験の聴取者：時間差の弁別閾（正答率75\,\%）は、帯域$150$--$1700$\,Hzの雑音で$9$\,\textmu s、$1$\,kHzの純音で$11$\,\textmu s、クリック音で$28$\,\textmu s。上限$0.6\ms$は章の計算、$0.22\,\text{m}/343\,\text{m/s}$ & \cite{KlumppEady1956} & 測定済み \\
8 & 鳥では遅延線と同時検出器によって時間差が場所に変わる~\eqref{eq:itd}（図~\ref{fig:itd}） & メンフクロウ：両耳からの軸索は同時検出器の核に反対側から入る。軸索に沿ったスパイクの到着時刻は深さとともに規則的に変わり、核を横切る遅延の幅（$700$\,\textmu mで$160$\,\textmu s）は両耳間遅延の範囲と一致する & \cite{Carr1990} & 測定済み \\
9 & 哺乳類では遅延の列は見つかっていない。同じ課題を、\nt{GLYC}ニューロンからの時間的に精密な抑制をもつ同時検出器が解く & スナネズミの内側上オリーブでのin vivo記録：応答は方向の地図をなさない。マイクロピペットでグリシンとその阻害薬を投与すると、時間的に精密な\emph{グリシン}抑制が遅延の動作範囲を決めていることがわかる。ここでの抑制は\nt{GABA}ではなく\nt{GLYC}による & \cite{Brand2002,Grothe2010} & 測定済み \\
10 & 強い帰還をもつ群はフリップフロップであり、自己維持活動は刺激の後も数秒続く（図~\ref{fig:bistable}） & 記憶した点への遅延眼球運動課題（遅延$1$--$6$\,s）中のサル前頭前皮質：課題に関係する$170$個のニューロンのうち$87$個が遅延中に発火率を変え、そのうち$79\,\%$では記憶した点の方向に依存する。この活動を二つの安定状態をもつ帰還が維持しているというのは理論である & \cite{Funahashi1989,Wang2001} & 測定済み \\
11 & 流入が$\approx0.7\tau$より長く続けばビットが書き込まれる（図~\ref{fig:recurrentgroup}(b)） & 方程式$\tau\,\dd r/\dd t=-r+f(wr+I)$を$w=1$、$I=0.6$でオイラー法により計算。\texttt{models/}にスクリプトはない。実験はない & 章内の計算 & モデル \\
12 & 記憶はほとんどスパイクなしにシナプスの促通に保存できる & 理論：終末の残留カルシウムが書き込んだ内容を約1秒保持する。根拠：フェレットの内側前頭前皮質（複数ニューロンの同時記録）には、後効果の長い促通性シナプスで結ばれた錐体ニューロンのサブネットワークがある。記憶保持中の「静かな」区間の観察の総説。皮質がいつどちらの方法を使うかは未解決の問題 & \cite{Mongillo2008,WangMarkram2006,Stokes2015} & 仮説 \\
13 & 記憶は疲労によって消去される。神経伝達物質の蓄えが枯渇する & 皮質錐体ニューロンのペア記録：高頻度のスパイクでシナプスの応答が低下し、低下の速さは放出確率で決まる。これこそが自己維持する群を消すことは直接には示されていない & \cite{Tsodyks1997} & 測定済み \\
14 & 群の連鎖はイベントで起動するタイマーであり、鳴禽ではニューロンが厳密に順番に発火する & 睡眠中と歌唱中のキンカチョウの高次発声中枢で、同定されたニューロンを記録：運動核に出力を送る各ニューロンは、歌の一つの正確な瞬間に短いバーストを1回出し、ニューロンは順に発火する & \cite{Hahnloser2002} & 測定済み \\
\bottomrule
\end{longtable}}
